\documentclass[11pt]{article}
\usepackage[margin=1in]{geometry}
\usepackage[T1]{fontenc}
\usepackage{lmodern}
\usepackage{amsmath,amssymb,amsthm}
\usepackage{graphicx}
\usepackage[hidelinks]{hyperref}
\usepackage{microtype}

\newtheorem{theorem}{Theorem}
\newcommand{\R}{\mathbb R}
\newcommand{\Ptwo}{\mathcal P_2}

\begin{document}

\begin{center}
{\Large Counterexample packet for arXiv:2311.15365}\par\vspace{5pt}
{\large Critical quadratic growth can destroy attainment}
\end{center}

\paragraph{Status.}
\textbf{counterexample\_likely\_valid}. This packet gives a negative answer
to the critical-\(p=2\) minimizer-existence question in Noboru Isobe,
\emph{A convergence result of a continuous model of deep learning via
\L{}ojasiewicz--Simon inequality}, arXiv:2311.15365v2~\cite{source}.

\paragraph{Source problem.}
The source considers, for
\(\eta\in L^2(0,1;\Ptwo(\R^m))\),
\begin{equation}\label{eq:sourceJ}
 J(\eta)=\int_{\mathsf X\times\mathsf Y}\ell\,d\mu_1^\eta
 +\frac{\lambda}{2}\int_0^1\int_{\R^m}|\theta|^2\,d\eta_t(\theta)\,dt,
\end{equation}
where \(\mu_t^\eta\) solves
\begin{equation}\label{eq:ce}
 \partial_t\mu_t^\eta+\operatorname{div}_x
 \bigl(\mu_t^\eta v_{\eta_t}\bigr)=0,
 \qquad
 v_{\eta_t}(x)=\int_{\R^m}v(x,\theta)\,d\eta_t(\theta).
\end{equation}
Its existence theorem assumes the subcritical bounds, for some \(p<2\),
\begin{align}\label{eq:growth}
 |v(x,\theta)|&\leq C|\theta|^p(1+|x|),\\
 |v(x_1,\theta)-v(x_2,\theta)|&\leq
 C|\theta|^p|x_1-x_2|.
\end{align}
Remark 7.1 on PDF page 23 says that ``the scenario when \(p=2\) in
Assumption 1.3 remains unsolved'' and asks for an example in which a minimizer
does not exist.

\begin{figure}[ht]
 \centering
 \includegraphics[width=\textwidth]{figures/open_problem_crop.png}
 \caption{The open question in Remark 7.1 of arXiv:2311.15365v2, PDF page 23.}
\end{figure}

\paragraph{Proof intuition.}
The regularization charges the second moment of the parameter distribution.
We choose a vector field whose ratio to this quadratic cost is strictly less
than one at every finite parameter but tends to one at infinity. A mass of
order \(M^{-2}\) placed at \(\theta=M\) then has a fixed quadratic cost and a
nonvanishing effect on the terminal state. The limiting effect-to-cost ratio
is never achieved by a finite parameter distribution, so the objective has a
minimizing sequence but no minimizer.

\begin{theorem}[Explicit nonattainment at \(p=2\)]\label{thm:main}
In the source problem, take
\[
 d=m=1,\qquad \mathsf X=\R,\qquad \mathsf Y=\{0\},\qquad
 \mu_0=\delta_{(0,0)},\qquad \lambda=2,
\]
and define
\begin{equation}\label{eq:data}
 v(x,\theta)=\frac{\theta^4}{1+\theta^2},
 \qquad
 \ell(x,0)=(x-1)^2.
\end{equation}
Then all assumptions in the source's minimizer-existence problem hold with
the critical exponent \(p=2\), but
\[
 \inf_{\eta\in L^2(0,1;\Ptwo(\R))}J(\eta)=\frac34
\]
and the infimum is not attained.
\end{theorem}

\begin{proof}
First check the hypotheses. The set \(\mathsf Y\) and the support of \(\mu_0\)
are compact, \(\lambda>0\), and \(\ell\) is continuous, nonnegative, and
\[
 \ell(x,0)=(x-1)^2\leq 2+2x^2.
\]
The vector field is continuous and independent of \(x\). Moreover,
\[
 0\leq v(x,\theta)=\theta^2\frac{\theta^2}{1+\theta^2}
 \leq \theta^2,
\]
whereas
\(v(x_1,\theta)-v(x_2,\theta)=0\). Thus both bounds in
\eqref{eq:growth} hold with \(p=2\) and \(C=1\).

Fix an arbitrary admissible \(\eta\), and put
\begin{equation}\label{eq:AR}
 R=\int_0^1\int_\R\theta^2\,d\eta_t(\theta)\,dt,
 \qquad
 A=\int_0^1\int_\R\frac{\theta^4}{1+\theta^2}
 \,d\eta_t(\theta)\,dt.
\end{equation}
Admissibility gives \(R<\infty\), and \(0\leq A\leq R\). Since the
averaged vector field in \eqref{eq:ce} is independent of \(x\), its flow is
translation by
\[
 x_t=\int_0^t\int_\R\frac{\theta^4}{1+\theta^2}
 \,d\eta_s(\theta)\,ds.
\]
Starting from \(\mu_0=\delta_{(0,0)}\), this gives
\(\mu_1^\eta=\delta_{(A,0)}\). Hence, because \(\lambda/2=1\),
\begin{equation}\label{eq:JAR}
 J(\eta)=(A-1)^2+R.
\end{equation}

There is a strict pointwise efficiency loss at every nonzero finite
parameter:
\begin{equation}\label{eq:strict}
 R-A=\int_0^1\int_\R\frac{\theta^2}{1+\theta^2}
 \,d\eta_t(\theta)\,dt.
\end{equation}
If \(R>0\), then \(\eta_t\) is not concentrated at zero for a set of times
of positive measure, and \eqref{eq:strict} is strictly positive. Thus
\(A<R\). If \(R=0\), then \(A=0\) and \(J(\eta)=1\). In the remaining case,
\begin{align*}
 J(\eta)
 &>(A-1)^2+A\\
 &=\left(A-\frac12\right)^2+\frac34
 \geq\frac34.
\end{align*}
Consequently every admissible \(\eta\) satisfies \(J(\eta)>3/4\).

It remains to show sharpness. Put \(r=1/2\), and for every \(M\geq1\) take
the constant-in-time curve
\begin{equation}\label{eq:seq}
 \eta_t^M=\left(1-\frac{r}{M^2}\right)\delta_0
 +\frac{r}{M^2}\delta_M,
 \qquad 0\leq t\leq1.
\end{equation}
It belongs to \(L^2(0,1;\Ptwo(\R))\), and its quantities in
\eqref{eq:AR} are
\[
 R_M=r=\frac12,
 \qquad
 A_M=\frac{r}{M^2}\frac{M^4}{1+M^2}
 =r\frac{M^2}{1+M^2}\longrightarrow r=\frac12.
\]
Therefore
\[
 J(\eta^M)=(A_M-1)^2+\frac12\longrightarrow
 \left(\frac12-1\right)^2+\frac12=\frac34.
\]
The lower bound and this sequence prove that the infimum is \(3/4\); its
strictness for every admissible \(\eta\) proves nonattainment.
\end{proof}

\paragraph{Verification notes.}
The proof uses no external theorem beyond the source's formulation of the
continuity-equation model. The flow calculation can also be checked directly
in distributions because the time-dependent velocity is spatially constant
and integrable. The strict step is exact: the integrand in
\eqref{eq:strict} vanishes only at \(\theta=0\). The included script evaluates
the exact formulas for sample \(M\); it is a smoke test, not part of the proof.

\paragraph{Why the critical exponent matters.}
The measures in \eqref{eq:seq} converge narrowly to \(\delta_0\), but their
second moments stay equal to \(1/2\). At every finite \(M\), velocity is
strictly smaller than quadratic cost; the missing efficiency appears only at
infinity. For \(p<2\), a bounded second moment controls this lower-order
effect uniformly at infinity. At \(p=2\), the effect and the regularizer have
the same homogeneity, permitting precisely this escape-of-mass mechanism.

\paragraph{Novelty check and limitations.}
On 2026-07-21, the run's registry, solution, attempt, and proof-gap indexes
were searched for arXiv:2311.15365, the exact title, ``continuous DNN,''
``\(p=2\),'' ``quadratic growth,'' and ``minimizer does not exist.'' Bounded
primary-source searches on arXiv used the exact id together with the same
phrases and close ``continuous model of deep learning'' variants. They found
the source paper, currently v2, but no later paper or revision resolving this
question. This supports candidate novelty but is not an exhaustive claim
about all published literature. The theorem addresses minimizer existence at
the critical growth exponent; it does not address the source's separate
questions about analyticity, curvature, or convergence of gradient flows.

\paragraph{Human-review recommendation.}
Review as a full negative answer to the critical-\(p=2\) existence question
and as the explicit nonattainment example requested in Remark 7.1. The two
points most worth auditing are the translation formula for \(\mu_1^\eta\) and
the implication \(R>0\Rightarrow R-A>0\).

\begin{thebibliography}{9}
\bibitem{source}
N. Isobe,
\emph{A convergence result of a continuous model of deep learning via
\L{}ojasiewicz--Simon inequality},
arXiv:2311.15365v2, 2024.
\end{thebibliography}

\end{document}
